\documentclass[11pt]{article}
\usepackage[margin=1in]{geometry}
\usepackage{amsmath,amssymb,amsthm}
\usepackage{graphicx}
\usepackage{microtype}
\usepackage{xcolor}
\usepackage{hyperref}
\hypersetup{colorlinks=true,linkcolor=blue!55!black,urlcolor=blue!55!black,citecolor=blue!55!black}
\newtheorem{theorem}{Theorem}
\newcommand{\Pp}{\mathbb P}
\newcommand{\gm}{\mathbin{\sharp}}
\newcommand{\gmt}{\mathbin{\sharp_t}}
\title{The printed geometric-mean conjecture already fails for scalars\\in arXiv:2401.00337}
\author{Candidate counterexample packet}
\date{11 August 2026}

\begin{document}
\maketitle

\begin{center}
\fbox{\parbox{0.91\linewidth}{\textbf{Status:} candidate counterexample likely valid. This is a full negative answer to Conjecture 3.7 exactly as printed. It is also a semantic warning: the likely intended $t=\tfrac12$ or $\sharp_t$ formulation is not settled here.}}
\end{center}

\section{Source conjecture}

The source defines the weighted geometric mean
\[
 A\gmt B=A^{1/2}(A^{-1/2}BA^{-1/2})^tA^{1/2},\qquad 0\leq t\leq1,
\]
and writes $A\gm B$ for the unweighted case $t=\tfrac12$.  Its final
Conjecture 3.7 asks, for $1\leq s\leq2$, whether
\begin{equation}\label{eq:printed}
 \left\|\sum_{i=1}^m (A_i^s\gm B_i^s)^r\right\|
 \leq
 \left\|\left[
 \left(\sum_iA_i\right)^{(1-t)srp/2}
 \left(\sum_iB_i\right)^{tsrp}
 \left(\sum_iA_i\right)^{(1-t)srp/2}
 \right]^{1/p}\right\|
\end{equation}
for $p>0$, $r\geq1$, $rp\geq1$, and every unitarily invariant norm~\cite{FH}.
The parameter $t$ appears on the right but is neither introduced nor
quantified in the conjecture; the left side contains $\gm$, not $\gmt$.

\begin{center}
\includegraphics[width=\textwidth]{figures/conjecture_crop.png}

\smallskip
\emph{The complete conjecture and open-case sentence in the source PDF, page 12.}
\end{center}

\section{Proof Intuition}

The displayed inequality is required for matrices of every size, so it must
survive the one-dimensional case.  There all noncommutativity disappears and
the unweighted mean on the left contributes equal powers of $A$ and $B$.
The right side instead contributes the unequal powers $1-t$ and $t$.  Their
ratio is a single power of $A/B$, which can be made larger than one unless
$t=\tfrac12$.  The counterexample therefore targets the exact mismatch in the
printed display; it does not rely on delicate matrix behavior.

\section{Counterexample and exact scalar obstruction}

\begin{theorem}\label{thm:counterexample}
Under the natural reading that $t\in[0,1]$ is inherited from the preceding
weighted-mean statements, \eqref{eq:printed} is false for every $s>0$ as
soon as $t\neq\tfrac12$.  In particular it is false in the stated open range
$1<s<2$.
\end{theorem}

\begin{proof}
Take $n=m=1$, so positive definite matrices are positive scalars and every
unitarily invariant norm is a positive multiple of the absolute value.  Put
$A_1=A>0$ and $B_1=B>0$.  Because
\[
 A^s\gm B^s=(A^sB^s)^{1/2}=(AB)^{s/2},
\]
the two sides of \eqref{eq:printed} reduce, after cancelling the norm's
constant factor, to
\[
 L=(AB)^{sr/2},\qquad R=A^{(1-t)sr}B^{tsr}.
\]
Consequently
\[
 \frac{L}{R}=\left(\frac AB\right)^{sr(t-1/2)}.
\]
If $t>\tfrac12$, choose $A>B$; if $t<\tfrac12$, choose $A<B$.  In either
case $L/R>1$.  Hence scalar validity for all positive $A,B$ forces
$t=\tfrac12$.

For a concrete counterexample in the open range, choose
\[
 A=4,\quad B=1,\quad s=\frac32,\quad r=p=1,\quad t=1.
\]
All printed hypotheses are satisfied, while
\[
 L=(4\cdot1)^{3/4}=2\sqrt2>1=4^0 1^{3/2}=R.
\]
\end{proof}

The same computation shows that the printed endpoint assertions cannot hold
for arbitrary $t$: the defect is independent of $s$.  This internal
inconsistency is further evidence that the final display lost either a
subscript $t$ on the left or the specialization $t=\tfrac12$ on the right.

\section{What remains open after repairing the statement}

There are two coherent repairs:
\begin{enumerate}
 \item replace $A_i^s\gm B_i^s$ by $A_i^s\gmt B_i^s$ while retaining
 $t\in[0,1]$ on the right; or
 \item fix $t=\tfrac12$ on the right and retain the unweighted geometric mean
 on the left.
\end{enumerate}
For scalars, both repairs give equality.  The source's Theorem 3.6 is an
unweighted result for $s\geq2$, so the second repair is especially consistent
with the sentence immediately preceding the conjecture.  This packet makes
no claim about either noncommutative repaired problem.

As a guard against mistaking a typographical counterexample for a substantive
one, we ran a focused numerical upgrade search on the second repair.  Random
positive-definite $2\times2$ and $3\times3$ matrices (up to four summands),
several values of $s\in(1,2)$, $r$, $p$, Schatten norms, and Ky Fan norms
gave 280,000 matrix configurations and 700,000 Ky Fan tests, with no
violation; the largest ratio was
$0.9999999999798115$.  A differential-evolution maximization over four
noncommuting real $2\times2$ positive-definite matrices, across 50 parameter
regimes, likewise produced no ratio above $1$ (the largest was
$0.9999994503007108$ among nontrivial reported runs).  These computations are
not a proof, but they support keeping the repaired conjecture open rather than
claiming more than the scalar diagnosis establishes.

\section{Verification and novelty scope}

The accompanying SymPy script evaluates the displayed scalar example exactly
and verifies $L-R=2\sqrt2-1>0$.  No floating-point arithmetic is used in that
verification.  The two additional scripts reproduce the random and optimized
checks via
\begin{verbatim}
conda run --no-capture-output -n sandbox python \
  code/search_repaired_conjecture.py
conda run --no-capture-output -n sandbox python \
  code/optimize_repaired_conjecture.py
\end{verbatim}

A bounded search through 11 August 2026 covered the run's registry, solution,
attempt, and proof-gap indexes; the exact conjecture text; the arXiv identifier
and title with \emph{Conjecture 3.7}, \emph{counterexample}, and
\emph{$t$-geometric mean}; and later papers surfaced by those searches.  It
found no later source explicitly correcting or resolving Conjecture 3.7.
Novelty confidence is moderate pending specialist review.  The literal scalar
calculation is independent of novelty.

\section{Conclusion}

Conjecture 3.7 is false exactly as displayed: scalar validity already forces
$t=\tfrac12$, whereas its right side retains a free $t$ and its left side is
unweighted.  The example $A=4$, $B=1$, $s=3/2$, $r=p=t=1$ is a full explicit
counterexample.  Human review should treat this as a semantic warning and
decide which repaired noncommutative conjecture was intended.

\begin{thebibliography}{9}
\bibitem{FH}
S. Freewan and M. Hayajneh,
\emph{Norm inequalities involving geometric means},
arXiv:2401.00337 (2024).
\url{https://arxiv.org/abs/2401.00337}.
\end{thebibliography}

\end{document}
